\chapter{Relation de Chasles}

\begin{itemize}
\item Vérifier la relation $\overline{AC} = 
\overline{AB}+\overline{BC}$, les points $A$, $B$ et $C$ appartenant à un axe $(xy)$, dans les cas suivants : 
\begin{itemize}
\it Les abscisses des points $A$, $B$, $C$ sont $0$,
$+3$, $+5$.
\it Les abscisses des points $A$, $B$, $C$ sont $0$,
$-2$, $-4$.
\it Les abscisses des points $A$, $B$, $C$ sont $0$,
$+7$, $+2$.
\it Les abscisses des points $A$, $B$, $C$ sont $0$,
$-5$, $-3$.
\it Les abscisses des points $A$, $B$, $C$ sont $0$,
$+2$, $-3$.
\it Les abscisses des points $A$, $B$, $C$ sont $0$,
$-3$, $+2$.
\end{itemize}
\item Soit un axe $(xy)$, $O$ l'origine des abscisses, vérifier la relation $\overline{AB} = 
\overline{OB}-\overline{OA}$ et déterminer la longueur $AB$ dans les cas suivants : \begin{itemize}
\it Les abscisses de $A$ et $B$ sont $+3$, $+5$.
\it Les abscisses de $A$ et $B$ sont $-2$, $-4$.
\it Les abscisses de $A$ et $B$ sont $+7$, $+2$.
\it Les abscisses de $A$ et $B$ sont $-5$, $-3$.
\it Les abscisses de $A$ et $B$ sont $+2$, $-3$.
\it Les abscisses de $A$ et $B$ sont $-3$, $+2$.
\end{itemize}
\item Soient $A$ et $B$ deux points d'un axe, $M$ le milieu de $[AB]$; vérifier la relation $\displaystyle\overline{OM} = \frac{\overline{OA}+\overline{OB}}2$
dans les cas suivants : \begin{itemize}
\it Les abscisses de $A$ et $B$ sont $+5$, $+9$.
\it Les abscisses de $A$ et $B$ sont $-5$, $+9$.
\it Les abscisses de $A$ et $B$ sont $-5$, $-9$.
\it Les abscisses de $A$ et $B$ sont $-9$, $+5$.
\end{itemize}
\it Soient sur un axe $(xy)$ les points $A$, $B$, $C$, $D$ d'abscisses $+2$, $+4$, $-1$, $-3$. Vérifier la relation $\overline{AD} = \overline{AB} + \overline{BC}
+\overline{CD}$ ; en est-il ainsi, quelle que soit la disposition des points $A$, $B$, $C$, et $D$ ?
\it Soient quatre points $A$, $B$, $C$, $D$ sur un axe. Démontrer la relation : $$\overline{AB} + \overline{BC} + \overline{CD}
+\overline{DA}=0.$$
\it Calculer le nombre d'années qui s'est écoulé entre les dates suivantes : $51$ ans avant J.-C. et 800 ans après J.-C. Généraliser avec les années $a$ et $b$. Tenir compte du fait qu'il n'y a pas d'année zéro.
\it Soient $A$ et $B$ d'abscisses $+4$ et $-3$. Déterminer les abscisses des points $C$ et $D$ qui partagent $[AB]$ en trois parties égales. Généraliser en désignant par $a$ l'abscisse de $A$ et par $b$ celle de $B$. 
\it Soient $A$, $B$, $C$, et $D$ quatre points d'un axe dont les abscisses sont $+2$, $-3$, $-1$ et $+4$. Vérifier la relation : \[ \overline{DA}.\overline{BC}+\overline{DB}.\overline{CA}+\overline{DC}.\overline{AB}=0.\]
En est-il toujours ainsi lorsqu'on désigne par $a$, $b$, $c$, $d$ les abscisses des points $A$, $B$, $C$ et $D$.
\end{itemize}